% Options for packages loaded elsewhere
\PassOptionsToPackage{unicode}{hyperref}
\PassOptionsToPackage{hyphens}{url}
%
\documentclass[
]{article}
\usepackage{amsmath,amssymb}
\usepackage{lmodern}
\usepackage{ifxetex,ifluatex}
\ifnum 0\ifxetex 1\fi\ifluatex 1\fi=0 % if pdftex
  \usepackage[T1]{fontenc}
  \usepackage[utf8]{inputenc}
  \usepackage{textcomp} % provide euro and other symbols
\else % if luatex or xetex
  \usepackage{unicode-math}
  \defaultfontfeatures{Scale=MatchLowercase}
  \defaultfontfeatures[\rmfamily]{Ligatures=TeX,Scale=1}
\fi
% Use upquote if available, for straight quotes in verbatim environments
\IfFileExists{upquote.sty}{\usepackage{upquote}}{}
\IfFileExists{microtype.sty}{% use microtype if available
  \usepackage[]{microtype}
  \UseMicrotypeSet[protrusion]{basicmath} % disable protrusion for tt fonts
}{}
\makeatletter
\@ifundefined{KOMAClassName}{% if non-KOMA class
  \IfFileExists{parskip.sty}{%
    \usepackage{parskip}
  }{% else
    \setlength{\parindent}{0pt}
    \setlength{\parskip}{6pt plus 2pt minus 1pt}}
}{% if KOMA class
  \KOMAoptions{parskip=half}}
\makeatother
\usepackage{xcolor}
\IfFileExists{xurl.sty}{\usepackage{xurl}}{} % add URL line breaks if available
\IfFileExists{bookmark.sty}{\usepackage{bookmark}}{\usepackage{hyperref}}
\hypersetup{
  pdftitle={Does environmental quality influence where people live?},
  hidelinks,
  pdfcreator={LaTeX via pandoc}}
\urlstyle{same} % disable monospaced font for URLs
\usepackage[margin=1in]{geometry}
\usepackage{graphicx}
\makeatletter
\def\maxwidth{\ifdim\Gin@nat@width>\linewidth\linewidth\else\Gin@nat@width\fi}
\def\maxheight{\ifdim\Gin@nat@height>\textheight\textheight\else\Gin@nat@height\fi}
\makeatother
% Scale images if necessary, so that they will not overflow the page
% margins by default, and it is still possible to overwrite the defaults
% using explicit options in \includegraphics[width, height, ...]{}
\setkeys{Gin}{width=\maxwidth,height=\maxheight,keepaspectratio}
% Set default figure placement to htbp
\makeatletter
\def\fps@figure{htbp}
\makeatother
\setlength{\emergencystretch}{3em} % prevent overfull lines
\providecommand{\tightlist}{%
  \setlength{\itemsep}{0pt}\setlength{\parskip}{0pt}}
\setcounter{secnumdepth}{-\maxdimen} % remove section numbering
\usepackage{booktabs}
\usepackage{longtable}
\usepackage{array}
\usepackage{multirow}
\usepackage{wrapfig}
\usepackage{float}
\usepackage{colortbl}
\usepackage{pdflscape}
\usepackage{tabu}
\usepackage{threeparttable}
\usepackage{threeparttablex}
\usepackage[normalem]{ulem}
\usepackage{makecell}
\usepackage{xcolor}
\ifluatex
  \usepackage{selnolig}  % disable illegal ligatures
\fi

\title{Does environmental quality influence where people live?}
\author{true}
\date{12-02-2021}

\begin{document}
\maketitle

{
\setcounter{tocdepth}{2}
\tableofcontents
}
\begin{verbatim}
## Rows: 3140 Columns: 16
\end{verbatim}

\begin{verbatim}
## -- Column specification --------------------------------------------------------
## Delimiter: ","
## chr (8): stfips, state_code, county_code, state_name, state, county_name, ru...
## dbl (8): cat_rucc, EQI, air_EQI, water_EQI, land_EQI, built_EQI, sociodemogr...
\end{verbatim}

\begin{verbatim}
## 
## i Use `spec()` to retrieve the full column specification for this data.
## i Specify the column types or set `show_col_types = FALSE` to quiet this message.
\end{verbatim}

This post discusses a statistical analysis used to answer the question:
does environmental quality influence where people live in the United
States?

\hypertarget{background}{%
\subsection{Background}\label{background}}

The Environmental Quality Index (EQI), developed by the U.S.
Environmental Protection Agency (EPA) provides a county level snapshop
of environmental conditions throughout the country. EPA first released
EQIs for the period 2000-2005 and updated these indexes for 2006-2010.
This statistical evaluation focuses on the 2006-2010 EQI. The purpose of
the EQI is to use (1) as an indicator of ambient conditions/exposure in
environmental health and modeling and (2) as a covariate to adjust for
ambient conditions in environmental models (EPA 2020). Previous studies
have used the EQI to evaluate relationships between environmental
quality and public health outcomes such as preterm birth, mortality,
cancer incidence, asthma, obesity, infant mortality

The EQI is developed from 5 environmental domains each with identified
environmental constructs as show in Table No.~1. Each county has an
overall environmental index and a domain specific index. Indexes were
also stratified by rural-urban continuum codes (RUCCs) for counties
classified as metropolitan urbanized, non-metro urbanized, less
urbanized, and thinly populated.

\begin{table}

\caption{\label{tab:unnamed-chunk-3}EQI Environmental Domains and Constructs}
\centering
\begin{tabular}[t]{>{}l||l}
\hline
Domain & Constructs\\
\hline
\textbf{air} & criteria air pollutants and hazardous air pollutants\\
\hline
\textbf{water} & overall water quality, general water contamination, domestic use, atmospheric deposition, drought, chemical contamination, and drinking water quality\\
\hline
\textbf{land} & agriculture, pesticides, facilities, radon, and mining activity\\
\hline
\textbf{built} & roads, highway/road safety, commuting behavior, housing environment, walkability, and green space\\
\hline
\textbf{sociodemographic} & crime, socioeconomic, political character, and creative class representation\\
\hline
\end{tabular}
\end{table}

\hypertarget{limitations}{%
\subsubsection{Limitations}\label{limitations}}

While the EQI can identify counties with higher environmental burdens,
it may not identify environmental injustices at local community level.
The EQI can not quantify environmental exposure for individuals and
reflects only outside environmental conditions, not indoor conditions.
The EQI can be used identify locations for future research, but is not
intended for regulatory purposes or as a diagnostics tool. Due to
changes in methodology and datasets, the 2000-2005 and 2006-2010 EQIs
should not be directly compared.

\hypertarget{the-data}{%
\subsection{The Data}\label{the-data}}

\hypertarget{environmental-quality-index}{%
\subsubsection{Environmental Quality
Index}\label{environmental-quality-index}}

To develop the EQI, variables were identified from available data to
represent each environmental domain and assessed for collinearity so
redundant variables could be excluded. Then, variables were standardized
based on geographic space or on a per capita rate, as appropriate. Data
gaps were evaluated to distinguish between missing data and meaningful
zeros. Where applicable, spatial kriging was used to interpolate values
when data was not available for all counties. For example, crime likely
occurs in all counties despite centralized reporting areas. Finally, the
normality of each variable was assessed and transformations such as
log-transformations, performed as needed.

Principal component analysis (PCA) was used to aggregate variables into
domain specific indexes. The domain indexes were then aggregated into
overall indexes for each county. A result of this method is that each
domain does not equally influence the overall EQI value for a given
county. The EQI is developed to be normally distributed with mean=0 and
standard deviation=1.

Higher EQI values correspond with worse environmental quality. Lower
(more negative) EQI values correspond with better environmental quality

\hypertarget{census-population-data}{%
\subsubsection{Census Population Data}\label{census-population-data}}

County level population data was obtained from the U.S. Census Bureau's
county intercensal datasets for 2000-2010. Percent population change was
calculated for 2006-2010 then winsorized to remove outliers above the
99.9th percentile (ie. counties with population change above 28.5\%).

\hypertarget{statistical-methods}{%
\subsection{Statistical Methods}\label{statistical-methods}}

Kalawao County, Hawaii had the lowest EQI value (highest environmental
quality) and greatest decrease in population (-17.4\%). Falls Church,
Virginia had the highest EQI (lowest environmental quality) and a
population change of 15.4\% which falls is the 99th percentile. The
summary statistics in table 2 suggest that population growth is
correlated with worse environmental characteristics.

\hypertarget{hypothesis-testing}{%
\subsubsection{Hypothesis testing}\label{hypothesis-testing}}

A different means test was completed to determine if mean environmental
quality was statistically different for counties that experienced
positive vs negative population change between 2006-2010. The figure
below shows a histogram of EQI values for counties with negative and
positive population change.

\includegraphics{the-epas-environmental-quality-index-eqi-vs-population-change_files/figure-latex/unnamed-chunk-5-1.pdf}

\begin{verbatim}
## Saving 5 x 3 in image
\end{verbatim}

\textbf{null\_hypothesis:} There is no difference in mean EQI for
counties with positive and negative population change.

\[H_{0}: \mu_{posPopChange} - \mu_{negPopChange} = 0\]

\textbf{alternative hypothesis:} There is a difference in EQI for
counties with positive and negative population change.

\[H_{A}: \mu_{posPopChange} - \mu_{negPopChange} \neq 0\]

\[\text{point estimate} = \mu_{posPopChange} - \mu_{negPopChange} =0.131 - -0.245 = 0.376\]

The standard error for a difference in means is defined as:

\[SE = \sqrt{\frac{s_1^2}{n_1} + \frac{s^2_2}{n_2}} = 0.036\]

and the z-score for hypothesis testing is:

\[z = \frac{\text{point estimate - null}}{SE} = \frac{0.376 - 0}{0.036} = 10.443\]

p-value, probability of getting a point estimate at least as extreme as
calculated if the null hypothesis were true:

\[p\text {- value }=\operatorname{Pr}(Z<-|z| \text { or } Z>|z|)=2 * \operatorname{Pr}(Z>|z|) = \ensuremath{1.5809578\times 10^{-25}}\]

Since the p-value \textless{} 0.001 we reject the null that there is no
difference in EQI for counties with positive population change versus
negative population change. There is a statistically significant
difference (at the 0.1\% significance level) in EQI across the two
population change groups. The 95\% confidence interval ranges from 0.31
to 0.45. This means that there is a 95\% chance that this interval
includes the true difference in mean EQI between counties with positive
and negative percent population change.

\hypertarget{linear-regression}{%
\subsubsection{Linear Regression}\label{linear-regression}}

Linear regression was used to model the relationship between
environmental quality and population change using the overall EQI value
and each domain specific EQI to determine if a particular domain was a
stronger predictor of population change.

\[\text{percent population change}_i=\beta_{0}+\beta_{1} \cdot EQI_i + \varepsilon_i\]

\includegraphics{the-epas-environmental-quality-index-eqi-vs-population-change_files/figure-latex/unnamed-chunk-11-1.pdf}

\begin{verbatim}
## Saving 5 x 3 in image
\end{verbatim}

First, hypothesis testing was used to test whether the slope coefficient
for the \% population change rate is equal to zero or not.

\textbf{null\_hypothesis:} The slope coefficient is equal to zero

\[H_{0}: \beta_{1} = 0\] \textbf{alternative hypothesis:} The slope
coefficient is \textbf{NOT} equal to zero

\[H_{A}: \beta_{1} \neq 0\]

point estimate, \(\beta_1\) = 0.991

SE = 0.077

the z-score for hypothesis test is:

\[z = \frac{\text{point estimate - null}}{SE} = \frac{{0.991 - 0}}{0.077} = 12.817\]

p-value: the probability of getting a point estimate at least as extreme
as calculates above if the null hypothesis were true:

\[p\text {- value }=\operatorname{Pr}(Z<-|z| \text { or } Z>|z|)=2 * \operatorname{Pr}(Z>|z|) = \ensuremath{1.0849376\times 10^{-36}}\]

Since p-value for the slope coefficient \textless{} 0.001, we reject the
null hypothesis that EQI has no influence on population change at the
0.01\% level. There is a statistically significant relationship between
EQI and \% population change and the coefficient is significantly
different from zero. The 95\% confidence interval for the slope
coefficient ranges from 0.84 to 1.143. This means that there is a 95\%
chance that this interval includes the true county level rate of change
for percent population change for each one unit change in EQI.

\hypertarget{domain-specific-linear-models}{%
\subsubsection{Domain Specific Linear
Models}\label{domain-specific-linear-models}}

\[\text{percent population change}_i=\beta_{0\text_domain}+\beta_{1\text_domain} \cdot EQI_i\text_domain + \varepsilon_i\]

\includegraphics{the-epas-environmental-quality-index-eqi-vs-population-change_files/figure-latex/unnamed-chunk-16-1.pdf}

\begin{verbatim}
## Saving 6 x 4 in image
\end{verbatim}

While a partly manual method was used above, the table below present
coefficients for domain specific models generated with statistical
functions in R. The numbers in {[}brackets{]} are the 95\% confidence
intervals for each estimated coefficient.

\begin{verbatim}
## Rows: 16 Columns: 7
\end{verbatim}

\begin{verbatim}
## -- Column specification --------------------------------------------------------
## Delimiter: ","
## chr (7): coefficient, overall EQI, air model, water model, land model, built...
\end{verbatim}

\begin{verbatim}
## 
## i Use `spec()` to retrieve the full column specification for this data.
## i Specify the column types or set `show_col_types = FALSE` to quiet this message.
\end{verbatim}

\begin{table}
\centering
\begin{tabular}{>{}l|l|l|l|l|l|l}
\hline
\textbf{coefficient} & \textbf{overall EQI} & \textbf{air model} & \textbf{water model} & \textbf{land model} & \textbf{built model} & \textbf{sociodem model}\\
\hline
\textbf{\textbf{Intercept}} & \textbf{1.868 ***} & \textbf{1.869 ***} & \textbf{1.868 ***} & \textbf{1.868 ***} & \textbf{1.867 ***} & \textbf{1.868 ***}\\
\hline
\textbf{} & [1.716, 2.019] & [1.717, 2.020] & [1.714, 2.023] & [1.713, 2.024] & [1.713, 2.021] & [1.718, 2.019]\\
\hline
\textbf{\textbf{EQI}} & \textbf{0.991 ***} & \textbf{} & \textbf{} & \textbf{} & \textbf{} & \textbf{}\\
\hline
\textbf{} & [0.840, 1.143] &  &  &  &  & \\
\hline
\textbf{\textbf{air EQI}} & \textbf{} & \textbf{0.998 ***} & \textbf{} & \textbf{} & \textbf{} & \textbf{}\\
\hline
\textbf{} &  & [0.847, 1.150] &  &  &  & \\
\hline
\textbf{\textbf{water EQI}} & \textbf{} & \textbf{} & \textbf{0.485 ***} & \textbf{} & \textbf{} & \textbf{}\\
\hline
\textbf{} &  &  & [0.331, 0.640] &  &  & \\
\hline
\textbf{\textbf{land EQI}} & \textbf{} & \textbf{} & \textbf{} & \textbf{-0.314 ***} & \textbf{} & \textbf{}\\
\hline
\textbf{} &  &  &  & [-0.469, -0.158] &  & \\
\hline
\textbf{\textbf{built EQI}} & \textbf{} & \textbf{} & \textbf{} & \textbf{} & \textbf{0.635 ***} & \textbf{}\\
\hline
\textbf{} &  &  &  &  & [0.481, 0.790] & \\
\hline
\textbf{\textbf{sociodem EQI}} & \textbf{} & \textbf{} & \textbf{} & \textbf{} & \textbf{} & \textbf{1.068 ***}\\
\hline
\textbf{} &  &  &  &  &  & [0.917, 1.219]\\
\hline
\textbf{n} & 3140 & 3140 & 3140 & 3140 & 3140 & 3140\\
\hline
\textbf{R2} & 0.050 & 0.051 & 0.012 & 0.005 & 0.020 & 0.058\\
\hline
\end{tabular}
\end{table}

The figure below provides a visual comparison of each model result. The
bold portion of the line represents the 90\% confidence interval and the
full line represents the 95\% confidence interval for each estimate.

\begin{verbatim}
## Loading required namespace: broom.mixed
## Loading required namespace: broom.mixed
## Loading required namespace: broom.mixed
## Loading required namespace: broom.mixed
## Loading required namespace: broom.mixed
## Loading required namespace: broom.mixed
## Loading required namespace: broom.mixed
## Loading required namespace: broom.mixed
## Loading required namespace: broom.mixed
## Loading required namespace: broom.mixed
## Loading required namespace: broom.mixed
## Loading required namespace: broom.mixed
\end{verbatim}

\includegraphics{the-epas-environmental-quality-index-eqi-vs-population-change_files/figure-latex/unnamed-chunk-18-1.pdf}
The p-value on the slope coefficient was \textless{} 0.05 for all domain
specific linear models which indicates a statistically significant
relationship. Based on the \(R^2\) values and slope coefficients, the
air and sociodemographic domains account for most of the overall
relationship between population change and EQI. The \(R^2\) terms
represents the variance in percent population change that can be
explained by EQI. For the overall EQI value, 5\% of the variance in
percent population change is explained by environmental conditions. The
sociademographic EQI explains 5.8\% of the variance in population change
while the air EQI explains 5.1\%. All domains except land are positively
correlated with population change. Since higher EQI values indicate
poorer environmental quality, these models show that population change
is greatest in counties with worse environmental conditions.

\hypertarget{conclusions}{%
\subsection{Conclusions}\label{conclusions}}

The identified relationships between population change and environmental
quality are noteworthy for their public health and environmental justice
implications. Increased population trends in areas with worse
environmental conditions could result increased incidences of cancer,
asthma, obesity, and infant mortality. While this project did not
evaluate economic variables, locations with higher environmental quality
could also have higher living costs which drive people to move to more
affordable places. If economic factors contribute to population growth
in counties with poor environmental quality, then this could negatively
affect the health of vulnerable populations and perpetuate social
inequalities. Further analysis could evaluate trends in mean household
income to determine if there is growing income inequality between
counties with better and worse environmental quality. Economic variables
and other factors influencing demographic shifts from rural areas cities
may be stronger predictors of population change than environmental
quality. In the future, EPA plans to create a census tract level EQI.
This would facilitate a higher resolution statistical evaluation of the
relationship between environmental conditions and population change.

\textbf{Data availability:}\\
\href{https://edg.epa.gov/data/public/ORD/CPHEA/}{EPA Datasets and files
from the EQI county data from 2006-2010}\\
\href{https://www.census.gov/data/tables/time-series/demo/popest/intercensal-2000-2010-counties.html}{U.S.
Census Bureau County Intercensal Tables: 2000-2010}

\hypertarget{references}{%
\subsection{References}\label{references}}

\begin{enumerate}
\def\labelenumi{\arabic{enumi}.}
\item
  U.S. EPA. Environmental Quality Index - Technical Report (2006-2010)
  (Final, 2020). U.S. Environmental Protection Agency, Washington, DC,
  EPA/600/R-20/367, 2020.
\item
  U.S. Census Bureau. County Intercensal Datasets: 2000-2010.
  \url{https://www.census.gov/data/datasets/time-series/demo/popest/intercensal-2000-2010-counties.html}
\end{enumerate}

\end{document}
